\id{ҒТАМР 50.47.31, 67.15.35,87.53.13}{}

\begin{header}
\swa{К.~Акишев, Л.~Акишева, Ж.~Нуртай}{ТЕХНОГЕНДІК ҚАЛДЫҚТАРДЫҢ ГРАНУЛОМЕТРИЯСЫН АВТОМАТТАНДЫРЫЛҒАН ТАЛДАУ}

\tsp{1}К.~Акишев\envelope,
\tsp{2}Л.~Акишева,
\tsp{1}Ж.~Нуртай
\end{header}

\begin{affil}
\tsp{1}К. Құлажанов атындағы Қазақ технология және бизнес университеті, Астана, Қазақстан,

\tsp{2}Назарбаев Зияткерлік мектебі, Астана, Қазақстан

\corrauthor{Корреспондент-автор: akmail04cx@mail.ru}
\end{affil}

Зерттеудің өзектілігі құрылыс саласында пайдалану кезінде металлургиялық
және энергетикалық өндірістің техногендік қалдықтарының гранулометриялық
құрамын жедел және сенімді бақылау қажеттілігіне байланысты. Дәстүрлі
әдістер, атап айтқанда Елек талдауы, жоғары еңбек сыйымдылығымен және
жаппай және автоматтандырылған бақылау жағдайында шектеулі қолданылуымен
сипатталады, бұл сандық талдау әдістерін енгізудің орындылығын
анықтайды.

Жұмыста компьютерлік көру және машиналық оқыту технологияларын қолдана
отырып, sem кескіндерін талдау негізінде қож және күл түйіршіктерінің
гранулометриялық құрамын анықтаудың кешенді әдісі ұсынылған. Бұл әдіс
1240 кескінді алдын-ала өңдеу, U-net моделін қолдана отырып бөлшектерді
сегментациялау, YOLOv5 негізіндегі объектілерді анықтау, геометриялық
белгілерді алу және Random forest, xgboost және SVM Алгоритмдер
ансамблін қолдана отырып фракцияларды жіктеу қадамдарын қамтиды.
Нәтижелерді верификациялау алынған үлестірімдерді МЕМСТ 5459.2-20 және
ASTM C136 сәйкес эталондық Елек талдауымен салыстыру арқылы орындалды.

Жұмыстың ғылыми жаңалығы айқын морфологиялық гетерогенділігі бар
техногендік қалдықтарды талдау үшін "компьютерлік көру - Машиналық
оқыту" гибридті архитектурасын қолдану болып табылады. Сандық нәтижелер
орташа абсолютті қателік 6-7\% болған кезде жіктеу дәлдігінің 95\% - ға
дейін жеткенін көрсетті, ал сапалық талдау күрделі пішіндегі және
бөлшектердің тығыз орналасуындағы әдістің тұрақтылығын растады.
Зерттеудің практикалық маңыздылығы қайталама шикізаттың сапасын
бақылауды автоматтандыру және талдау уақытын қысқарту мүмкіндігі болып
табылады. Әрі қарайғы зерттеулердің болашағы онлайн-түсірілімге көшумен,
өндірістік желілермен интеграцияланумен және техногендік материалдардың
деректер жиынтығын кеңейтумен байланысты.

{\bfseries Түйін сөздер}: машиналық оқыту, модельдер, гранулометриялық
құрам, техногендік қалдықтар, машиналық көру, алгоритм.

\begin{header}
АВТОМАТИЗИРОВАННЫЙ АНАЛИЗ ГРАНУЛОМЕТРИИ ТЕХНОГЕННЫХ ОТХОДОВ

\tsp{1}К.~Акишев\envelope,
\tsp{2}Л.~Акишева,
\tsp{1}Ж.~Нуртай
\end{header}

\begin{affil}
\tsp{1}Казахский университет технологии и бизнеса им. К. Кулажанова, Астана, Казахстан,

\tsp{2}Назарбаев интеллектуальная школа, Астана, Казахстан,

e-mail: akmail04cx@mail.ru
\end{affil}

Актуальность исследования обусловлена необходимостью оперативного и
достоверного контроля гранулометрического состава техногенных отходов
металлургического и энергетического производства при их использовании в
строительной отрасли. Традиционные методы, в частности ситовой анализ,
характеризуются высокой трудоёмкостью и ограниченной применимостью в
условиях массового и автоматизированного контроля, что обуславливает
целесообразность внедрения цифровых методов анализа.

В работе предложен комплексный метод определения гранулометрического
состава гранул шлака и золы на основе анализа SEM-изображений с
использованием технологий компьютерного зрения и машинного обучения.
Метод включает этапы предварительной обработки 1240 изображений,
сегментации частиц с применением модели U-Net, детекции объектов на
основе YOLOv5, извлечения геометрических признаков и классификации
фракций с использованием ансамбля алгоритмов Random Forest, XGBoost и
SVM. Верификация результатов выполнена путём сопоставления полученных
распределений с эталонным ситовым анализом согласно ГОСТ 5459.2-20 и
ASTM C136.

Научная новизна работы заключается в применении гибридной архитектуры
«компьютерное зрение - машинное обучение» для анализа техногенных
отходов с выраженной морфологической неоднородностью. Количественные
результаты показали достижение точности классификации до 95 \% при
средней абсолютной ошибке 6 - 7 \%, а качественный анализ подтвердил
устойчивость метода при сложной форме и плотной укладке частиц.
Практическая значимость исследования заключается в возможности
автоматизации контроля качества вторичного сырья и сокращения времени
анализа. Перспективы дальнейших исследований связаны с переходом к
онлайн-съёмке, интеграцией с производственными линиями и расширением
датасета техногенных материалов.

{\bfseries Ключевые слова}: машинное обучение, модели, гранулометрический
состав, техногенные отходы, машинное зрение, алгоритм.

\begin{header}
AUTOMATED ANALYSIS OF TECHNOGENIC WASTE GRANULOMETRY

\tsp{1}K.~Akishev\envelope,
\tsp{2}L.~Akisheva,
\tsp{1}Zh.Nurtai
\end{header}

\begin{affil}
\tsp{1}K. Kulazhanov Kazakh University of Technology and Business, Astana, Kazakhstan,

\tsp{2}Nazarbayev Intellectual School, Astana, Kazakhstan,

e-mail: akmail04cx@mail.ru
\end{affil}

The relevance of the study is due to the need for prompt and reliable
control of the granulometric composition of man-made waste from
metallurgical and energy production when used in the construction
industry. Traditional methods, in particular sieve analysis, are
characterized by high labor intensity and limited applicability in
conditions of mass and automated control, which makes it advisable to
introduce digital analysis methods.

The paper proposes a comprehensive method for determining the
granulometric composition of slag and ash granules based on the analysis
of SEM images using computer vision and machine learning technologies.
The method includes the stages of 1240 image preprocessing, particle
segmentation using the U-Net model, YOLOv5-based object detection,
geometric feature extraction, and fraction classification using an
ensemble of Random Forest, XGBoost, and SVM algorithms. Verification of
the results was performed by comparing the obtained distributions with
the reference sieve analysis according to GOST 5459.2-20 and the ASTM
C136 standard.

The scientific novelty of the work lies in the application of the hybrid
architecture "computer vision - machine learning" for the analysis of
man-made waste with pronounced morphological heterogeneity. The
quantitative results showed the achievement of classification accuracy
of up to 95\% with an average absolute error of 6 -7\%, and the
qualitative analysis confirmed the stability of the method in complex
shapes and dense packing of particles. The practical significance of the
research lies in the possibility of automating the quality control of
secondary raw materials and reducing the analysis time. The prospects
for further research are related to the transition to online
photography, integration with production lines, and expansion of the
dataset of man-made materials.

{\bfseries Keywords}: machine learning, models, granulometric composition,
man-made waste, machine vision, algorithm.

\begin{multicols}{2}
{\bfseries Кіріспе.} Тау-кен металлургия өндірісінің техногендік қалдықтары
кен материалдарын өндіру, байыту және өңдеу кезінде түзілетін қайталама
ресурстардың едәуір көлемін білдіреді. Бұл қалдықтарға көбінесе мөлшері,
пішіні, тығыздығы және минералды құрамы бойынша ерекшеленетін ұсақ және
ірі фракциялық бөлшектер жатады. Гранулометриялық құрам техногендік
қалдықтарды одан әрі пайдалану мүмкіндігін айқындайтын негізгі
параметрлердің бірі болып табылады: құрылыста, 3D-басып шығаруда,
металлургиялық өңдеуде және өнеркәсіптің басқа да салаларында құрам
үшін.

Елек талдауы және лазерлік дифракция сияқты гранулометриялық құрамды
анықтаудың дәстүрлі әдістері айтарлықтай уақыт пен материалдық
шығындарды, сондай-ақ үлгілерді қолмен өңдеуді қажет етеді. Деректер
көлемінің ұлғаюы және бөлшектердің параметрлерін жедел және дәл бағалау
қажеттілігі жағдайында заманауи цифрлық және интеллектуалды
технологияларды қолдануға қызығушылық артады.

Машиналық оқыту (ML) кескіндерді, спектрлерді және басқа бөлшектердің
сипаттамаларын талдауды автоматтандыруға кең мүмкіндіктер береді. Сандық
микроскоптар немесе жоғары ажыратымдылықтағы камералар арқылы алынған
жасанды материалдардың кескіндерін өңдеу үшін терең оқыту және
компьютерлік көру әдістерін қолдану әсіресе перспективалы.

Ғылыми зерттеулердің қазіргі жағдайын талдау әртүрлі материалдардың,
соның ішінде техногендік қалдықтардың гранулометриялық құрамын талдау
үшін машиналық оқытуды қолдануға қызығушылықтың тұрақты өсуін көрсетті.
Бұл шолу осы мәселені шешудің әртүрлі тәсілдерін жүзеге асыратын
жұмыстарды қарастырады.

Соңғы жылдары Машиналық оқыту және компьютерлік көру әдістері минералды
бөлшектердің гранулометриялық құрамы мен морфологиялық сипаттамаларын
анықтау міндеттерінде кеңінен дамыды. Бөлшектердің үлестірілуін өлшем
бойынша бағалауды автоматтандыруға сегменттеу мен детекцияның терең
нейрондық модельдеріне негізделген жұмыстар айтарлықтай үлес қосты.

Бірқатар зерттеулер instance segmentation-дің бөлшектердің мөлшерін
анықтаудағы тиімділігін көрсетеді. Сонымен {[}1{]} нейрондық желіні
сегментациялау арқылы минералды бөлшектердің PSD есептеу әдісі ұсынылды,
бұл зертханалық әдістермен салыстыруға болатын нәтижелерге қол жеткізді.
Ұқсас бағытты nie және басқалар дамытады {[}2,3{]} және басқалар, мұнда
терең сегменттеу модельдері (U-Net және оның модификациялары)
микроскопиялық кескіндерде минералды дәндерді бөлудің жоғары дәлдігін
қамтамасыз етеді. Алайда, бұл зерттеулер негізінен табиғи минералдарға
және бақыланатын зертханалық жағдайларға бағытталған, бұл олардың
құрылымында гетерогенділігі жоғары техногендік қалдықтарға тікелей
қолданылуын шектейді.

Көп деңгейлі және ансамбльдік архитектураларды қолданатын жұмыстар жеке
бағытты құрайды. Зерттеулерде {[}4,5{]} көп масштабты және ансамбльдік
тәсілдер бөлшектердің күрделі морфологиясында сегменттеу тұрақтылығын
арттыратыны көрсетілген. Осыған қарамастан, бұл жұмыстарда модельдердің
архитектуралық аспектілеріне баса назар аударылады, статистикалық
тұрақтылық, кросс-валидация және анықтамалық талдау әдістерімен
салыстыру мәселелері қайталама болып қала береді.

Шолу басылымдары {[}6,7{]} минералды кескіндерді сегменттеу
мәселелерінде терең оқытудың заманауи әдістерін жүйелейді және
шешілмеген негізгі мәселелерді атап көрсетеді: деректер жиынтығының
сипаттамаларының ашықтығы, нәтижелердің нашар түсіндірілуі және
өнеркәсіптік жағдайларға бағытталған зерттеулердің шектеулі саны. Бұл
тұжырымдар дәлдікті, тұрақтылықты және практикалық қолдануды
біріктіретін жан-жақты зерттеулердің қажеттілігін растайды.

Онлайн бақылау мен өнеркәсіптік енгізуге бағытталған жұмыстар {[}8,9{]}
зерттеулерде ұсынылған, мұнда PSD талдауының жеделдігі мен
автоматтандырылуының маңыздылығы көрсетілген. Сонымен қатар, мұндай
жұмыстардың көпшілігінде әртүрлі модельдерді егжей-тегжейлі салыстыру
жоқ немесе нәтижелердің тұрақтылығының статистикалық негіздемесі
берілмейді.

Сондай-ақ, мамандандырылған архитектуралар мен доменді бейімдеуді
қолданатын зерттеулерді атап өткен жөн {[}10,11{]}, олар әртүрлі кескін
түрлері арасында модельдерді тасымалдау әлеуетін көрсетеді. Дегенмен,
бұл тәсілдер айтарлықтай есептеу ресурстарын қажет етеді және әдетте
жоғары мамандандырылған сценарийлерге (КТ деректері, шағылысқан жарық,
SEM кескіндері) бағытталған.

Жеке топ бөлшектердің пішіні мен морфологиясын талдауға арналған
жұмыстардан тұрады {[}12-14{]}, сондай-ақ өнеркәсіптік қалдықтарды қайта
өңдеуге байланысты зерттеулер {[}15{]}. Бұл басылымдар техногендік
материалдардың сапасын бағалау кезінде деректерді статистикалық өңдеу
мен инженерлік түсіндірудің маңыздылығын көрсетеді.

Терең оқыту әдістерінің кең таралуына дейін компьютерлік көруді
пайдалана отырып, бөлшектердің өлшемдері мен таралуын талдауға арналған
жұмыстарда суреттерге негізделген гранулометриялық құрамды анықтау
әдістерінің дамуына айтарлықтай үлес қосылды. Зерттеуде {[}16{]}
бөлшектерді оқшаулау әдістерін, перспективалық бұрмалануларды түзетуді
және өлшемді бөлуді статистикалық өңдеуді қоса алғанда, тау жыныстарының
фрагментациясын талдау үшін цифрлық кескінді өңдеудің негізгі
принциптері құрылды. Бұл жұмыс image-based гранулометрия саласындағы
іргелі жұмыстардың бірі ретінде қарастырылады және бүгінгі күнге дейін
әдістемелік негіз ретінде қолданылады.

Image-based тәсілдерінің одан әрі дамуы {[}17{]} еңбектерінде көрінеді,
онда үйінділердегі бөлшектердің көрінуін жіктеу міндеті қарастырылады
және қабаттасу мен төсеу тығыздығының өлшемді бағалау дәлдігіне әсері
көрсетілген. Бұл нәтижелер бөлшектердің қабаттасу геометриясы
суреттерден гранулометриялық құрамды талдауда маңызды фактор болып
табылатындығын және нәтижелерді түсіндіру кезінде ескерілуі керек екенін
көрсетті.

Kemeny {[}18{]} өз зерттеуінде морфологиялық кескін сегментациясын
қолдана отырып, темір кені түйіршіктерінің өлшемдерін өлшеуге арналған
өнеркәсіптік машиналық көру жүйесін ұсынады. Авторлар өндіріс жағдайында
image-based әдістерін практикалық қолдану мүмкіндігін көрсетті және
алынған өлшемді бөлуді стандартты өлшеу нәтижелерімен салыстырды. Бұл
бағыттың дамуы {[}19{]} және басқалардың зерттеуінде жалғасты мұнда
конвейерлік желілердегі әктас материалдары үшін бөлшектердің массалық
таралуын бағалау жүйесі ұсынылған, бұл өнеркәсіпте image-based
гранулометриясын қолдану аясын едәуір кеңейтті.

Осылайша, дереккөздер {[}16-19{]} осы зерттеуде қолданылатын Машиналық
оқыту мен нейрондық желі модельдерін қолдану тәсілдерін қоса алғанда,
заманауи әдістерге негізделген гранулометриялық құрамды анықтаудың
классикалық және инженерлік негізделген image-based негізін құрайды.

Осылайша, әдебиеттерді талдау қолданыстағы зерттеулер негізінен
минералды бөлшектерді сегментациялауға немесе терең нейрондық желілердің
архитектуралық аспектілеріне бағытталғанын көрсетеді, бұл ретте
модельдерді кешенді салыстыру, статистикалық тұрақтылық және техногендік
қалдықтарды талдаудың анықтамалық әдістерімен салыстыру мәселесі
жеткілікті түрде ашылмаған. Бұл жұмыс осы олқылықтың орнын толтыруға
бағытталған:

- техногендік қалдықтардың жаңғыртылатын деректер жиынтығын пайдалану;

- гибридті тәсілді қолдану (сегменттеу/анықтау + ML-жіктеу);

- кросс-валидация мен сенімділік интервалдарын қолдана отырып
модельдерді статистикалық салыстыру;

- нәтижелерді өнеркәсіптік эталон ретінде Елек талдауымен салыстыру.

{\bfseries Материалдар мен әдістер.} \emph{Қож және күл түйіршіктеріне
арналған машиналық көру жүйесінің жұмыс кезеңдері.}

Техногендік материалдардың гранулометриялық құрамын анықтау үшін
сканерлеуші электронды микроскопия (SEM) әдісімен алынған қож және күл
түйіршіктерінің цифрлық кескіндерін талдауға негізделген машиналық көру
жүйесі пайдаланылды сур.1.
\end{multicols}

\fig[0.6\textwidth]{i/image20}[1 - сурет. Magus Микроскопы]

\begin{multicols}{2}
SEM кескіндері жоғары кеңістіктік ажыратымдылықпен сипатталады және
бөлшектердің нақты морфологиялық ерекшеліктерін, соның ішінде дұрыс емес
пішінді, бетінің кедір-бұдырын, ұсақ фракциялардың агломерациясын және
өлшемдердің кең ауқымын түсіруге мүмкіндік береді, бұл оларды құрылыс
мақсатындағы техногендік қалдықтарды зерттеудің өкілі етеді.

Бірінші кезеңде машиналық көру жүйесі үшін кіріс ретінде пайдаланылатын
бастапқы кескіндер (2, a-сурет) жасалады. Кескіндер жарықтық пен
контраст бойынша алдын-ала қалыпқа келтіріледі, бұл түсіру жағдайлары
мен материалдың гетерогенділігі өзгерген кезде талдаудың тұрақтылығын
қамтамасыз етеді.

Екінші кезеңде бөлшектерді сегментациялау жүзеге асырылады (2, b-сурет),
фонда жеке түйіршіктерді оқшаулауға бағытталған. Сегменттеу контурларды,
текстуралық және морфологиялық белгілерді талдау арқылы жүзеге
асырылады, нәтижесінде жеке бөлшектерге сәйкес келетін пиксельдік
маскалар пайда болады. Күл және қож материалдарына тән агломерацияланған
аймақтарды бөлуге ерекше назар аударылады, бұл ұсақ фракцияларды дұрыс
ескеруге мүмкіндік береді.

Үшінші кезең бөлшектерді анықтауды қамтиды (2-сурет, c), онда әрбір
сегменттелген түйіршік жеке объект ретінде қарастырылады. Детекция
кескін кеңістігіндегі бөлшектердің локализациясын қамтамасыз етеді,
олардың геометриялық шекараларын анықтауға және бөлшектердің қабаттасуы
мен тығыз орналасуына байланысты қателерді азайтуға мүмкіндік береді.

Соңғы кезеңде (2, d-сурет) бөлшектердің геометриялық белгілерін, оның
ішінде ауданы мен эквивалентті диаметрін алу жүзеге асырылады, олардың
негізінде гранулометриялық құрамның таралуы қалыптасады. Алынған
мәліметтер фракцияларды жіктеу үшін қолданылады және ұсынылған әдістің
верификациясы мен практикалық қолданылуын қамтамасыз ететін эталондық
Елек талдауының нәтижелерімен салыстырылады.

Ұсынылған қадамдар техногендік материалдардың гранулометриялық құрамын
автоматтандырылған анықтауға бағытталған SEM кескіндерін дәйекті өңдеуді
көрсетеді. Сегменттеуді қолдану күрделі пішінді бөлшектерді дұрыс
оқшаулауға және ұсақ фракцияларды ескеруге мүмкіндік береді, ал детекция
Объектілік талдау деңгейін қамтамасыз етеді және қабаттасудың әсерін
азайтады. Алынған геометриялық белгілерге сүйене отырып, бөлшектердің
мөлшерін бөлу қалыптасады, ол фракцияларды жіктеу және кейіннен електен
талдау нәтижелерімен салыстыру арқылы тексеру үшін қолданылады. Бұл
тәсіл нәтижелердің қайталануын және құрылыс индустриясындағы токсиндер
мен күлдің сапасын бақылау әдісін практикалық қолдану мүмкіндігін
қамтамасыз етеді.
\end{multicols}

\begin{figs}
\fig[0.3\textwidth]{i/image21}\hspace{-3cm}
\fig[0.3\textwidth]{i/image22}
\end{figs}
\vspace{-0.6cm}
\begin{figs}[2 - сурет. Әртүрлі мөлшердегі шлак пен күл түйіршіктері (A-d)]
\fig[0.3\textwidth]{i/image23}\hspace{-3cm}
\fig[0.3\textwidth]{i/image24}
\end{figs}

\begin{multicols}{2}
\emph{Сәулет және алгоритмдерді таңдау.} Суретте.3 модельдерді қосымша
оқыту үшін кері байланысы бар қож және күл түйіршіктерінің
гранулометриялық құрамын анықтау кезінде компьютерлік көру және
машиналық оқыту алгоритмдерінің өзара әрекеттесуінің Блок-схемасы
берілген.

Кезеңдер тізбегі:

- шлак пен күл түйіршіктерінің SEM суреттері;

- бөлшектердің нақты морфологиясы бар кіріс;

- алдын ала өңдеу;

- жарықтықты қалыпқа келтіру, шуды сүзу, контрастты күшейту;

- бөлшектерді сегментациялау-U-Net;

- пиксель деңгейі: жеке түйіршіктердің маскаларын қалыптастыру;

- бөлшектерді анықтау-YOLOv5;

- нысан деңгейі: bounding boxes, бөлшектерді локализациялау;

- белгілерді алу;

- ауданы, периметрі, эквивалентті диаметрі;

- фракциялардың жіктелуі-кездейсоқ Орман / XGBoost / SVM;

- өлшем аралықтары бойынша ансамбльдік жіктеу;

- гранулометриялық құрамы (PSD);

- ГОСТ (Елек талдауы) және ASTM c136 стандарттары бойынша бөлуді
қалыптастыру және тексеру.

Кері байланыс алынған гранулометриялық құрамды үлестіруді эталондық Елек
талдауының (ГОСТ, ASTM C136) {[}20,21{]} нәтижелерімен салыстыру арқылы
жүзеге асырылады. Жүйелі ауытқулар анықталған жағдайда жіктеуіштер мен
сегменттеу моделінің параметрлері нақтыланады, бұл жүйенің техногендік
материалдардың өзгеретін қасиеттеріне бейімделуін қамтамасыз етеді.
\end{multicols}

\fig[0.6\textwidth]{i/image25}[3-сурет. Модельдерді оқыту үшін кері байланысы бар қож және күл
түйіршіктерінің гранулометриялық құрамын анықтау кезінде компьютерлік
көру және машиналық оқыту алгоритмдерінің өзара әрекеттесуінің
блок-схемасы]

\begin{multicols}{2}
Бөлшектерді фракциялар бойынша жіктеу үшін шешім ағаштарының әдістері
мен ансамбльдік модельдерді қамтитын Машиналық оқыту алгоритмдерінің
ансамблі қолданылды. Бұл алгоритмдерді таңдау олардың шуға
төзімділігіне, сызықтық емес тәуелділіктермен жұмыс істеу қабілетіне
және нәтижелердің түсіндірілуіне байланысты.

Ансамбльдік тәсілді қолдану жіктеудің тұрақтылығын арттыруға және
техногендік материалдардың мөлшерін бөлуге тән шығарындылардың әсерін
азайтуға мүмкіндік берді.

Оқу үлгісі SEM кескіндерін өңдеу нәтижесінде алынған мәліметтер
негізінде қалыптасты. Әрбір бөлшек үшін үлгіге геометриялық белгілердің
мәндері және эталондық Елек талдауының нәтижелері бойынша анықталған
сәйкес фракциялық тиістілік енгізілді. Үлгі модельдердің сапасын дұрыс
бағалауды қамтамасыз ететін оқыту, валидация және тестілік кіші
үлгілерге бөлінді.

Модельдерді оқыту бекітілген параметрлерді және нәтижелердің қайталануын
қамтамасыз ететін деректерді бөлудің бірыңғай схемасын қолдану арқылы
жүзеге асырылды. Жіктеу сапасы стандартты көрсеткіштерді, соның ішінде
дәлдікті және фракциялар бойынша орташа абсолютті қатені пайдалана
отырып бағаланды.

Нәтижелерді верификациялау гранулометриялық құрамның алынған
үлестірімдерін нормативтік талаптарға сәйкес орындалған эталондық Елек
талдауының нәтижелерімен салыстыру арқылы жүзеге асырылды. Бұл тәсіл
ұсынылған әдістің дұрыстығын және оның құрылыс индустриясының
практикалық міндеттері үшін қолданылуын сандық бағалауға мүмкіндік
берді.

Машиналық оқыту әдістері жеткілікті кеңістіктік ажыратымдылықпен және
бөлшектердің орташа тығыздығымен кескіндерді талдауда ең жоғары
тиімділікті көрсетеді. Өте ұсақ фракциялар мен айқын қабаттасуы бар
кескіндер үшін жіктеу дәлдігі төмендейді, бұл ұсынылған тәсілдің
қолданылу шекарасын және оны одан әрі дамыту бағытын анықтайды.

Ұсынылған тәсілде Машиналық оқыту модельдері оқшауланбайды, бірақ
деректерді өңдеудің дәйекті каскады түрінде қолданылады, онда әр модель
қатаң анықталған мәселені шешеді және нәтижелерді келесі кезеңге береді.
Мұндай архитектура фракцияларды сегменттеу, анықтау және жіктеу
міндеттерін бөлуге, жүйенің тұрақтылығын арттыруға және техногендік
материалдардың гранулометриялық құрамын талдау кезінде нәтижелердің
қайталануын қамтамасыз етуге мүмкіндік береді. Эксперименттік деректер
жиынтық күні 1-кестеде келтірілген.
\end{multicols}

\tcap{1-кесте. Эксперименттік деректер жиынтығының сипаттамалары}
\begin{longtblr}[
  label = none,
  entry = none,
]{
  cells = {c},
  cells = {font = \small},
  row{1} = {font=\bfseries\small},
  hlines,
  vlines,
}
Параметр                    & Мағынасы                                  \\
Деректер түрі               & Техногендік қалдықтардың сандық бейнелері \\
Суреттер саны               & 1240                                      \\
Фракциялық сыныптардың саны & 6                                         \\
Фракция Диапазоны, мм       & 0–0,5; 0,5–1; 1–2; 2–5; 5–10; 10          \\
Кескін пішімі               & TIFF, PNG                                 \\
Кескін ажыратымдылығы       & 5–20 Мп                                   \\
Аннотация әдісі             & Қолмен белгілеу (LabelMe, CVAT)           \\
Белгілеуді тексеру          & 100 \%                                    \\
Деректерді бөлу             & Train/Validation/Test = 70/15/15          
\end{longtblr}

\begin{multicols}{2}
1- кестеде іріктеу көлемін, фракциялық сыныптар санын, кескін пішімін,
аннотация әдісін және деректерді оқыту, валидация және сынақ қосалқы
үлгілеріне бөлу схемасын қоса алғанда, эксперименттік деректер жиынының
негізгі параметрлері берілген. Кесте жүргізілген зерттеудің
репрезентативтілігі мен қайталануын бағалауға мүмкіндік береді.

\emph{U-net байламы → YOLOv5 (компьютерлік көру)}

U-net сегментациясы (пиксель деңгейі)

U-net моделі SEM кескіндеріндегі бөлшектерді сегменттеу үшін кескінді
өңдеудің бірінші кезеңінде қолданылады. Оның міндеті-қож және күл
түйіршіктері алып жатқан аймақтарға сәйкес келетін екілік маскаларды
қалыптастыру. Кодер-декодер архитектурасы мен skip қосылыстарының
арқасында U-Net ұсақ фракциялар мен агломерацияланған аймақтарды қоса
алғанда, күрделі пішінді бөлшектердің шекараларын дәл бөлуді қамтамасыз
етеді.

U-Net нәтижесі:

- бөлшектердің пиксельдік маскалары → анық бөлінген нысандары бар
тазартылған сурет.

Бұл жағдайда u-Net фракцияларды жіктемейтінін, бірақ объектіні талдау
үшін деректерді дайындайтынын түсіну керек.

\emph{YOLOv5-анықтау және бастапқы жіктеу (объект деңгейі)}

Келесі қадамда сегменттелген кескіндер YOLOv5 моделіне түседі, ол жеке
бөлшектерді анықтау мәселесін шешеді. YOLOv5 әр түйіршік үшін шектеу
тіктөртбұрыштарын (bounding boxes) құрайды және оның кеңістіктік
координаттарын анықтайды. Сонымен қатар, модель бөлшектерді Өлшем
диапазоны бойынша бастапқы жіктеу үшін қолданылады, бұл деректердің
Объектілік көрінісін автоматтандыруға мүмкіндік береді.

\emph{YOLOV5} нәтижесі:

- нысандар жиынтығы → координаттар, bounding boxes өлшемдері, бөлшектер
идентификаторы.

Осылайша, ауысу схема бойынша жүзеге асырылады:

- пикселдер → Нысандар.

\emph{YOLOV5 байламы → кездейсоқ Орман / XGBoost / SVM (кесте
деректері)}

Анықталғаннан кейін әр бөлшек сандық белгілер жиынтығымен сипатталады,
соның ішінде:

- ауданы, периметрі, эквивалентті диаметрі, bounding box жақтарының
қатынасы.

Бұл белгілер классикалық машиналық оқыту алгоритмдері қолданатын
деректердің кестелік көрінісін құрайды.

Random Forest және XGBoost алгоритмдері бөлшектерді фракциялық
интервалдар бойынша түпкілікті жіктеу үшін қолданылады. Оларды қолдану
шуға төзімділікке, сызықтық емес тәуелділіктермен жұмыс істеу қабілетіне
және нәтижелердің жоғары интерпретациясына байланысты. Бұл модельдердің
ансамбльдік сипаты техногендік материалдарға тән шығарындылардың әсерін
азайтуға мүмкіндік береді.

SVM моделі жіктеу тұрақтылығы мен нәтижелердің салыстырмалылығын тексеру
үшін бақылау алгоритмі ретінде қолданылады. SVM шығыстарын ансамбльдік
модельдермен салыстыру шешім қабылдаудың дұрыстығын растауға және
жіктеудің шекаралық жағдайларын анықтауға мүмкіндік береді.

Осылайша, ұсынылған машиналық оқыту жүйесі гранулометриялық құрамды
талдауға көп деңгейлі тәсілді жүзеге асырады: U-Net бөлшектердің пиксель
деңгейінде дәл сегменттелуін қамтамасыз етеді, YOLOv5 детекция мен
объектіні локализациялауды жүзеге асырады, ал Random Forest, XGBoost
және SVM алынған белгілерге негізделген фракцияларды жіктейді. Мұндай
модульдік құрылым талдаудың дәлдігін арттыруға, нәтижелердің қайталануын
қамтамасыз етуге және жүйені техногендік материалдардың әртүрлі
түрлеріне бейімдеуге мүмкіндік береді.

Алынған нәтижелердің дұрыстығы мен қайталануын қамтамасыз ету үшін осы
зерттеу шеңберінде модельдерді оқыту параметрлерін, қолданылатын
архитектуралардың ашықтығын және ұқсас деректерде эксперименттерді
қайталау мүмкіндігін қамтитын верификацияның кешенді әдістемесі іске
асырылды.

Біріншіден, модельдерді оқыту параметрлері эксперименттер барысында
тіркелді және өзгертілмеді, бұл гиперпараметрлерді субъективті таңдаудың
қорытынды нәтижелерге әсерін болдырмайды. Машиналық оқытудың барлық
модельдері үшін бірдей деректерді оқыту, валидация және тестілеу
үлгілеріне бөлу схемалары (70/15/15), random seed тұрақты мәндері,
сондай-ақ оқуды тоқтатудың бірыңғай критерийлері қолданылды. Бұл әртүрлі
модельдер арасындағы нәтижелердің салыстырмалылығын қамтамасыз етеді
және дәлдік пен тұрақтылық көрсеткіштеріндегі айырмашылықтарды дұрыс
түсіндіруге мүмкіндік береді.

Екіншіден, қолданылатын модельдердің архитектуралары ғылыми әдебиеттерде
ашық және кеңінен сипатталған, бұл зерттеу әдіснамасының ашықтығын
арттырады. Пайдаланылған нейрондық желі және ансамбльдік Алгоритмдер
(U-Net, YOLO тәсілдері, Random Forest, XGBoost, SVM) жалпыға қол жетімді
құрылымдар мен стандартты Машиналық оқыту кітапханалары негізінде жүзеге
асырылады. Бұл басқа зерттеушілерге жеке шешімдерді немесе жабық
бағдарламалық жасақтаманы қажет етпестен ұсынылған тәсілді қайталауға
мүмкіндік береді.

Үшіншіден, эксперименттердің қайталануы кескін сипаттамаларын, Аннотация
әдістемесін, алдын ала өңдеу процедураларын және қолданылатын бағалау
көрсеткіштерін қоса алғанда, кіріс деректері мен өңдеу кезеңдерін
егжей-тегжейлі сипаттау арқылы қамтамасыз етіледі.5 есе
кросс-валидацияны қолдану және сенімділік интервалдарын талдау алынған
нәтижелердің тұрақтылығын және модельдерді қайта оқытудың жоқтығын
қосымша растайды. Ұқсас деректерді пайдалану және сипатталған
эксперимент шарттарын сақтау кезінде нәтижелер салыстырмалы дәлдікпен
тәуелсіз зерттеулерде қайталануы мүмкін.

Осылайша, ұсынылған верификация әдісі нәтижелердің ғылыми негізділігін,
ашықтығын және қайталануын қамтамасыз етеді, бұл әдісті одан әрі дамыту
және оны Құрылыс және қайта өңдеу салаларының практикасына енгізу үшін
қажетті шарт болып табылады.

{\bfseries Нәтижелер және талқылау.} Қож және күл түйіршіктерінің SEM
кескіндеріне машиналық көру жүйесін қолдану нәтижелері зерттелетін
материалдардың күрделі морфологиясына қарамастан бөлшектерді сегменттеу
мен анықтаудың жоғары тиімділігін көрсетті. Бөлшектер дұрыс емес
пішінмен, беттің кедір-бұдырымен, сондай-ақ агломерацияланған
аймақтардың болуымен сипатталады, бұл дәстүрлі түрде автоматтандырылған
талдауды қиындатады. Дегенмен, сегменттеу және анықтау каскадын
пайдалану жеке түйіршіктерді дұрыс оқшаулауға және деректердің тұрақты
Объектілік көрінісін қамтамасыз етуге мүмкіндік берді.

4-суретте классикалық Елек талдауын және нейрондық және классикалық
модельдер ансамблін қоса алғанда, машиналық оқытудың әртүрлі әдістерін
пайдалана отырып, алынған тау-кен металлургиялық өндірісінің техногендік
қалдықтарының гранулометриялық құрамын бөлуді салыстыру ұсынылған.
\end{multicols}

\fig{i/image26}[4-сурет. Классикалық Елек талдауын және машиналық оқытудың
әртүрлі әдістерін, соның ішінде нейрондық және классикалық модельдер
ансамблін пайдалана отырып алынған техногендік қалдықтардың
гранулометриялық құрамының таралуын салыстыру]

\begin{multicols}{2}
Эталондық әдіс ретінде Елек талдауы қолданылды, оның нәтижелері нүктелі
сызықпен белгіленді. Салыстыру үшін бір графикте U-Net + YOLOv5, Random
Forest, XGBoost және SVM модельдері арқылы алынған қисықтар берілген.

График бөлшектердің бөлшек интервалдары бойынша салынған (0-0,5; 0,5--1;
1-2; 2-5; 5-10; \textgreater 10 мм), ординаттар осі бойынша әрбір
фракцияның массалық үлесі пайызбен кейінге қалдырылды. Бұл көрініс
бөлшектердің жеке өлшем диапазонындағы үлестірімдердің жалпы
консистенциясын да, модельдердің жергілікті ауытқуларын да дәл бағалауға
мүмкіндік береді.

Суреттен көріп отырғанымыздай, U-Net + YOLOv5 және XGBoost модельдері
барлық фракциялар бойынша ең үлкен сәйкестікті көрсетеді, әсіресе 1-2 мм
және 2-5 мм диапазонындағы електі талдау нәтижелерімен, олар техногендік
қалдықтарды өңдеуде ең технологиялық маңызды болып табылады. Random
Forest моделі жақын нәтижелерді көрсетеді, бірақ ол шағын және үлкен
фракцияларда біршама үлкен қателіктермен сипатталады. Сонымен қатар, SVM
моделі эталоннан ең үлкен ауытқуларды көрсетеді, әсіресе ұсақ бөлшектер
аймағында, бұл күрделі морфологиялық белгілерді талдау кезінде сызықтық
ядроның экспрессивті қабілетін шектеуге байланысты болуы мүмкін.

Осылайша, 1-сурет ансамбльдік және гибридті Машиналық оқыту әдістерін
қолдану классикалық Елек талдауына айтарлықтай аз уақыт шығындарымен
және процестерді автоматтандырудың жоғары деңгейімен жоғары сәйкестік
дәрежесіне қол жеткізуге мүмкіндік беретінін анық растайды.

5-сурет машиналық оқытудың ансамбльдік тәсілімен алынған техногендік
қалдықтардың фракцияларын жіктеудегі қателіктер матрицасын көрсетеді.
Абсцисса осі бойынша модель болжаған фракциялар кластары, ординат осі
бойынша -- анықтамалық Елек талдауы негізінде анықталған шынайы сыныптар
ұсынылған. Матрица ұяшықтарындағы мәндер әрбір бөлшек интервал үшін
дұрыс және қате жіктелген бөлшектердің пайызын көрсетеді.
\end{multicols}

\fig[0.5\textwidth]{i/image27}[5-сурет. Техногендік қалдықтардың фракцияларын жіктеудегі
қателіктер матрицасы]

\begin{multicols}{2}
Матрицаның диагональды элементтері дұрыс жіктелген бөлшектерге сәйкес
келеді және көптеген фракциялар үшін 88-94\% - дан асатын жоғары
мәндермен сипатталады. Бұл ұсынылған тәсілдің жоғары жіктеу дәлдігі мен
тұрақтылығын көрсетеді. Қателердің ең көп концентрациясы 0,5--1 мм және
1-2 мм, сондай-ақ 2-5 мм және 5-10 мм сияқты іргелес фракциялық
аралықтарда байқалады. Бұл әсер күтілетін және шекаралық өлшемдегі
бөлшектердің морфологиялық жақындығына, сондай-ақ кескіндердегі
бөлшектердің тығыз орналасуындағы оптикалық ажыратымдылықтың
шектеулеріне байланысты.

Айта кету керек, іс жүзінде жіктеудің өрескел қателіктері жоқ, онда ұсақ
фракциялардың бөлшектері қателесіп үлкендерге жіктеледі және керісінше.
Бұл сегменттеу және анықтау алгоритмдерінің дұрыс жұмыс істеуін,
сондай-ақ таңдалған белгі кеңістігінің кейінгі жіктеуге сәйкестігін
көрсетеді.

Осылайша, ұсынылған қате матрицасы қателіктердің негізгі үлесі
жергілікті екенін және гранулометриялық құрамның қорытынды таралуына
айтарлықтай әсер етпейтінін растайды. Бұл ұсынылған әдісті жедел бақылау
және техногендік материалдарды автоматтандырылған талдау міндеттерінде
практикалық қолдануға жарамды етеді.

6-суретте барлық зерттелетін машиналық оқыту үлгілері үшін Accuracy
көрсеткіші бойынша 5 есе кросс-валидация нәтижелерін ұсынады: U-Net +
YOLOv5, Random Forest, XGBoost және SVM. Абсцисса осі бойынша
кросс-валидация қатпарларының нөмірлері, ординат осі бойынша - жіктеу
дәлдігінің мәндері қойылады.

Бұл график бастапқы датасеттің әр түрлі кіші үлгілерінде оқыту кезінде
әр модельдің тұрақтылығы мен қайталануын бағалауға мүмкіндік береді.
Суреттен көріп отырғанымыздай, U-Net + yolov5 және XGBoost модельдері
қатпарлар арасындағы дәлдік мәндерінің ең аз таралуын көрсетеді, бұл
жоғары тұрақтылық пен Айқын қайта оқытудың жоқтығын көрсетеді. Random
Forest моделі практикалық қолдану үшін рұқсат етілген мәндер шегінде
қалып, қалыпты таралуды көрсетеді.
\end{multicols}

\fig{i/image26}[6 - сурет. Accuracy көрсеткіші бойынша 5 есе кросс-валидация
нәтижелері]

\fig[0.6\textwidth]{i/image28}[7-сурет. Машиналық оқытудың әр моделі үшін жіктеу дәлдігінің
орташа мәні (Accuracy)]

\begin{multicols}{2}
SVM моделі нәтижелердің ең үлкен өзгергіштігімен сипатталады, бұл
берілген әдістің оқу үлгісін таңдауға сезімталдығын көрсетеді және оның
техногендік қалдықтардың күрделі көп класты таралуын талдау үшін
неғұрлым шектеулі қолданылуын растайды.

Осылайша, 3-суретте көрсетілген кросс-валидация нәтижелері
гранулометриялық талдау тапсырмалары үшін ансамбльдік және гибридтік
модельдерді қолданудың орындылығын растайды және алынған нәтижелерге
қосымша статистикалық негіздеме ретінде қызмет етеді.

7-сурет 5 есе кросс-валидация нәтижелері бойынша есептелген 95\%
сенімділік аралықтарын көрсете отырып, әрбір Машиналық оқыту моделі үшін
жіктеу дәлдігінің орташа мәндерін (Accuracy) көрсетеді. Сенімділік
аралықтарын пайдалану алынған көрсеткіштердің статистикалық сенімділігін
және нәтижелердің белгісіздік дәрежесін бағалауға мүмкіндік береді.

Суреттен көріп отырғанымыздай, Accuracy-дің ең жоғары орташа мәні
XGBoost моделін көрсетеді, ал сенімділік аралығы минималды енге ие, бұл
нәтижелердің жоғары тұрақтылығы мен қайталануын көрсетеді. U-Net +
YOLOv5 моделі сонымен қатар терең сегменттеу мен объектілерді анықтау
комбинациясының тиімділігін растайтын тар сенімділік аралығымен жоғары
дәлдікті көрсетеді.

Random Forest моделі оқу деректерінің вариацияларына орташа
сезімталдықты көрсететін біршама кеңірек сенімділік интервалымен
сипатталады. SVM моделі үшін ең төменгі орташа дәлдік және ең кең
сенімділік аралығы байқалады, бұл күрделі гранулометриялық
үлестірімдерді талдау кезінде берілген тәсілдің тұрақтылығының
төмендігін көрсетеді.

Осылайша, 4-сурет бір классификаторлармен салыстырғанда ансамбльдік
Машиналық оқыту әдістерінің артықшылықтарын статистикалық негізделген
растауды қамтамасыз етеді және оларды техногендік қалдықтарды
автоматтандырылған талдау міндеттерінде қолданудың орындылығын
көрсетеді.

8-суретте 5 есе кросс-валидация арқылы алынған әрбір машиналық оқыту
моделі үшін Accuracy мәндерін бөлудің boxplot диаграммаларын ұсынады.
Визуализацияның бұл түрі медианалық мәндерді, квартильаралық диапазонды,
шығарындылардың болуын және жіктеу дәлдігінің жалпы таралуын бір уақытта
бағалауға мүмкіндік береді.
\end{multicols}

\fig[0.7\textwidth]{i/image29}[8 - сурет. Boxplot-әр Машиналық оқыту моделі үшін Accuracy
мәндерін бөлу диаграммалары]

\begin{multicols}{2}
Ұсынылған деректерден xgboost және U-Net + YOLOv5 модельдері тар
квартильаралық интервалмен және медиана мен орташаның жақын мәндерімен
сипатталатынын көруге болады, бұл нәтижелердің жоғары тұрақтылығын және
оқу үлгісін таңдауға минималды сезімталдықты көрсетеді. Random Forest
моделі біршама үлкен дисперсияны көрсетеді, бірақ инженерлік және
өндірістік тапсырмалар үшін қолайлы мәндер шегінде қалады.

SVM моделі үшін ең кең квартильаралық диапазон, сондай-ақ Медиананың
төменгі Accuracy мәндеріне ауысуы байқалады, бұл техногендік қалдықтар
фракцияларының көп класты классификациясы мәселесін шешуде осы әдістің
шектеулі тиімділігін растайды.

Осылайша, суретте көрсетілген boxplot талдауы 8 модельдердің сапасын
кешенді статистикалық бағалауды қамтамасыз ете отырып және осы зерттеу
шеңберінде ансамбльдік тәсілді таңдаудың негізділігін растай отырып,
орташа мәндер мен сенімділік интервалдарының нәтижелерін толықтырады.

2-кестеде Accuracy, Recall, Precision, F1-score және орташа абсолютті
қате (MAE) көрсеткіштері бойынша U-Net + Yolov5, Random Forest, XGBoost
және SVM үлгілерінің салыстырмалы сәйкестігі ұсынылған. Кесте ең тиімді
модельді таңдауды негіздеу үшін қолданылады.
\end{multicols}

\tcap{2-кесте. Сапаның негізгі көрсеткіштері бойынша модельдерді салыстыру}
\begin{longtblr}[
  label = none,
  entry = none,
]{
  cells = {c},
  cells = {font = \small},
  row{1} = {font=\bfseries\small},
  row{4} = {font=\bfseries\small},
  hlines,
  vlines,
}
Модель         & Дәлдік & Қайта шақыру & Дәлдігі & F1-score & MAE, \% \\
U-Net + YOLOv5 & 0,94   & 0,92         & 0,94    & 0,93     & 6,8     \\
Random Forest  & 0,93   & 0,91         & 0,92    & 0,92     & 7,2     \\
XGBoost        & 0,95   & 0,93         & 0,95    & 0,94     & 6,3     \\
SVM            & 0,89   & 0,87         & 0,88    & 0,87     & 8,1     
\end{longtblr}

3-кестеде жеке кросс-валидация қатпарларында машиналық оқытудың әрбір
моделі үшін алынған жіктеу дәлдігінің (Accuracy) мәндері, сондай-ақ
орташа мәндер мен стандартты ауытқулар берілген. Кесте модельдердің
тұрақтылығы мен тұрақтылығын сипаттайды.

\tcap{3-кесте. есе кросс-валидация нәтижелері (Дәлдік)}
\begin{longtblr}[
  label = none,
  entry = none,
]{
  cells = {c},
  cells = {font = \small},
  row{1} = {font=\bfseries\small},
  hlines,
  vlines,
}
Fold   & U-Net + YOLOv5 & Random Forest & XGBoost & SVM   \\
Fold 1 & 0,93           & 0,92          & 0,94    & 0,87  \\
Fold 2 & 0,94           & 0,93          & 0,95    & 0,88  \\
Fold 3 & 0,95           & 0,94          & 0,96    & 0,89  \\
Fold 4 & 0,94           & 0,93          & 0,95    & 0,90  \\
Fold 5 & 0,93           & 0,93          & 0,95    & 0,89  \\
Mean   & 0,94           & 0,93          & 0,95    & 0,89  \\
Std    & 0,008          & 0,007         & 0,007   & 0,011 
\end{longtblr}

4-кестеде деректердің шуға төзімділігі, іріктеу сезімталдығы, инференция
жылдамдығы және масштабталуы бойынша модельдерді сапалы салыстыру
жүргізілді. Кесте өнеркәсіптік енгізу үшін модельдердің инженерлік және
қолданбалы жарамдылығын көрсетеді.

\tcap{4-кесте. Тұрақтылық пен есептеу шығындары бойынша әдістерді салыстыру}
\begin{longtblr}[
  label = none,
  entry = none,
]{
  cells = {c},
  cells = {font = \small},
  row{1} = {font=\bfseries\small},
  hlines,
  vlines,
}
Критерий                          & U-Net + YOLOv5 & Random Forest & XGBoost & SVM      \\
Шуға төзімділік                   & Жоғары         & Орташа        & Жоғары  & Төмен    \\
Үлгіге сезімталдық                & Төмен          & Орташа        & Төмен   & Жоғары   \\
Инференс жылдамдығы               & Орташа         & Жоғары        & Орташа  & Жоғары   \\
Масштабтау                        & Жоғары         & Жоғары        & Жоғары  & Орташа   \\
Өнеркәсіптік енгізуге жарамдылығы & иә             & иә            & иә      & Шектеулі 
\end{longtblr}

5-кестеде классикалық Елек талдауы және дәлдік, талдау уақыты,
автоматтандыру деңгейі және адам факторының әсері бойынша ұсынылған ML
тәсілі салыстырылады. Кесте автоматтандырылған әдістің практикалық
артықшылықтарын көрсетеді.

\tcap{5-кесте. Классикалық және ML тәсілдерін салыстыру}
\begin{longtblr}[
  label = none,
  entry = none,
]{
  cells = {c},
  cells = {font = \small},
  row{1} = {font=\bfseries\small},
  hlines,
  vlines,
}
Көрсеткіш               & Електі талдау & ML-тәсіл  \\
Орташа дәлдік           & Эталон        & 93–95\%   \\
1 үлгіні талдау белбеуі & 15–20 мин     & 3–4 мин   \\
Автоматизация           & жок           & Толық     \\
Адам факторы            & Жоғары        & Минималды \\
Масштабтау              & Шектеулі      & Жоғары    
\end{longtblr}

9-сурет техногендік қалдықтар фракцияларының жіктелу дәлдігінің оқыту
процесінде қолданылатын бастапқы кескіндердің сапасына және машиналық
оқыту модельдерінің инференсіне тәуелділігін көрсетеді.

\fig{i/image30}[9-сурет. Техногендік қалдықтар фракцияларының жіктелу дәлдігінің
бастапқы кескіндердің сапасына тәуелділігі]

\begin{multicols}{2}
Бұл контекстегі кескіндердің сапасы кеңістіктік ажыратымдылықты,
жарықтандырудың біркелкілігін, шу деңгейін және бұлыңғырлық дәрежесін
қамтитын параметрлер жиынтығын білдіреді. Графикте САПАНЫҢ үш тән
деңгейі көрсетілген: жоғары, орташа және төмен, зертханалық және
өндірістік сценарийлердегі нақты түсірілім жағдайларын көрсетеді.

Суреттен көріп отырғанымыздай, суреттердің жоғары сапасымен модельдер
бөлшектердің шекараларын нақты визуализациялауға және сегменттеу және
анықтау алгоритмдерінің дұрыс жұмыс істеуіне байланысты максималды
дәлдік мәндерін көрсетеді. Орташа сапаға көшу кезінде бөлшектердің
контурларының ішінара жоғалуына және жергілікті жіктеу қателіктерінің
өсуіне байланысты дәлдіктің орташа төмендеуі байқалады. Accuracy-тің ең
маңызды құлдырауы суреттердің сапасы төмен болған кезде жазылады, мұнда
Шу мен бұлыңғырлықтың әсері фракцияларды жіктеу кезеңіне каскадты түрде
берілетін сегменттеу қателіктеріне әкеледі.

Бұл сурет суреттердің сапасы компьютерлік көру алгоритмдерінің дұрыс
жұмыс істеуі үшін маңызды фактор болып табылатындығын көрнекі растау
ретінде қызмет етеді. Ол ұсынылған тәсілді өнеркәсіптік енгізу кезінде
суретке түсіру шарттарын реттеу, кескінді алдын ала өңдеу әдістерін
қолдану және жабдықты калибрлеу қажеттілігін атап көрсетеді. Осылайша,
9-сурет әдістің негізгі шектеулерінің бірін негіздеу үшін қолданылады
және көрсетілген дәлдік көрсеткіштеріне қандай жағдайда қол
жеткізілетінін көрсетеді.

10-суретте толықтық көрсеткішінің (Recall) талданатын кескіндердегі
бөлшектердің қабаттасу тығыздығына тәуелділігі көрсетілген. Төсеу
тығыздығы бөлшектердің қабаттасу дәрежесін және олардың визуалды
бөлінуінің күрделілігін анықтайтын маңызды технологиялық параметр болып
табылады. Зерттеу барысында тығыздықтың үш тән деңгейі қарастырылды:
төмен, орташа және жоғары, әртүрлі үлгілерді дайындау сценарийлеріне
сәйкес келеді.
\end{multicols}

\fig{i/image31}[10 - сурет. Толықтық көрсеткішінің (Recall) талданатын
кескіндердегі бөлшектердің қабаттасу тығыздығына тәуелділігі]

\begin{multicols}{2}
График төмен және орташа төсеу тығыздығында машиналық оқыту модельдері
жоғары Recall мәндерін көрсететінін көрсетеді, бұл көптеген бөлшектердің
дұрыс анықталғанын және олардың тиісті фракциялық сыныптарға
жатқызылғанын көрсетеді. Төсеу тығыздығы жоғары болған кезде
бөлшектердің контурларының ішінара бірігуіне және жетіспейтін объектілер
санының өсуіне байланысты толықтықтың айтарлықтай төмендеуі байқалады.
Бұл әсер әсіресе визуалды айырмашылықтар минималды болатын іргелес
фракциялық Интервалдардың бөлшектері үшін айқын көрінеді.

10-сурет ұсынылған әдістің үлгіні дайындау шарттарына сезімталдығын
көрсету үшін қолданылады және төсеу тығыздығы автоматтандырылған
талдаудың дәлдігін шектейтін факторлардың бірі екенін растайды. Сонымен
қатар, нәтижелер төсеу тығыздығы жоғары болса да, recall төмендеуі
бақыланатын болып қала беретінін және соңғы гранулометриялық
үлестірімнің сыни бұрмалануына әкелмейтінін көрсетеді.

Осылайша, ұсынылған график үлгілерді дайындау процедурасын стандарттау
немесе ұсақ фракциялардың жоғары концентрациясы бар материалдарды талдау
кезінде бөлшектерді бөлудің қосымша әдістерін қолдану қажеттілігін
негіздейді.

11-сурет бөлшектердің әртүрлі өлшем диапазондары бойынша жіктеудің
орташа абсолютті қателігінің (MAE) таралуын көрсетеді: ұсақ, орташа және
үлкен фракциялар. Бұл график қандай фракциялық интервалдар
гранулометриялық құрамды анықтаудағы жалпы қателікке үлкен үлес
қосатынын бағалауға мүмкіндік береді.
\end{multicols}

\fig{i/image32}[11-сурет. Бөлшектердің әртүрлі өлшем диапазондары бойынша
жіктеудің орташа абсолютті қателігінің таралуы (MAE): ұсақ, орташа және
үлкен фракциялар]

\begin{multicols}{2}
11-суреттен көрініп тұрғандай, MAE-дің ең үлкен мәндері өлшемдері
қолданылатын оптикалық жүйенің ажыратымдылық шегіне жақындаған ұсақ
фракцияларға тән. Мұндай бөлшектер үшін шудың, агломерацияның және
пиксельді іріктеу қателіктерінің әсері артады, бұл дұрыс сегменттеу мен
жіктеуді қиындатады. Орташа фракциялар үшін қате айтарлықтай төмен, бұл
бөлшектердің өлшемдері мен кескін ажыратымдылығының оңтайлы қатынасына
байланысты. Ірі фракциялар MAE минималды мәндерімен сипатталады, өйткені
олардың контурлары айқын және компьютерлік көру алгоритмдерімен оңай
анықталады.

Бұл сурет өте кішкентай фракцияларды талдау кезінде әдісті шектеуді
сандық негіздеу үшін қолданылады және ML тәсілі мен классикалық Елек
талдауы арасындағы негізгі сәйкессіздіктер жергілікті екенін растайды.
Сонымен қатар, кішігірім фракциялар үшін қателіктердің жоғарылауы
гранулометриялық құрамның интегралды таралуына айтарлықтай әсер
етпейтінін және әдістің практикалық құндылығын төмендетпейтінін атап
өткен жөн.

11-сурет сонымен қатар жоғары ажыратымдылықты түсіру әдістері мен
бейімделген сегменттеу алгоритмдерін қолдану арқылы ұсақ бөлшектерді
талдаудың дәлдігін арттыруға бағытталған әрі қарайғы зерттеу бағыттарын
анықтауға негіз болады.

6-кесте тау-кен металлургия өндірісінің техногендік қалдықтарының
гранулометриялық құрамын анықтау үшін қолданылатын машиналық оқытудың
зерттелетін модельдерінің қорытынды рейтингін ұсынады. Рейтинг жіктеу
дәлдігін (Accuracy), толықтығын (Recall), фракцияларды анықтаудағы
орташа абсолютті қатені (MAE), кросс-валидация нәтижелерінің
тұрақтылығын, кіріс деректерінің шуға сезімталдығын, есептеу шығындарын
және өнеркәсіптік жағдайларда ықтимал қолдануды қоса алғанда, сандық
және сапалық көрсеткіштерді кешенді талдау негізінде орындалады.
\end{multicols}

\tcap{6-кесте. Модельдердің қорытынды рейтингі}
\begin{longtblr}[
  label = none,
  entry = none,
]{
  cells = {c},
  cells = {font = \small},
  row{1} = {font=\bfseries\small},
  column{1} = {font=\bfseries\small},
  cell{2}{2} = {font=\bfseries\small},
  hlines,
  vlines,
}
Ранг & Модель         & Негіздеме                                             \\
1    & XGBoost        & Ең жақсы дәлдік, минималды қателік, жоғары тұрақтылық \\
2    & U-Net + YOLOv5 & Дәл сегменттеу және анықтау, тұрақтылық               \\
3    & Random Forest  & Жақсы жалпылау, қарапайымдылық                        \\
4    & SVM            & Күрделі фракциялар үшін шектеулі тиімділік            
\end{longtblr}

\begin{multicols}{2}
Соңғы рейтингте ең жоғары Accuracy және F1-score мәндерін, MAE минималды
мәнін және 5 есе кросс-валидациядағы нәтижелердің ең аз таралуын
көрсеткен XGBoost моделі маңызды орын алады. Тар сенімділік аралықтары
мен көрсеткіштердің тұрақтылығы осы модельдің жоғары жалпылау қабілетін
көрсетеді, бұл бөлшектердің өзгермелі морфологиясы бар гетерогенді
техногендік материалдарды талдауда өте маңызды.

Екінші орынды u-net + yolov5 гибридті тәсілі алады, ол кескіндердегі
бөлшектерді сегменттеу мен анықтаудың жоғары дәлдігін қамтамасыз етеді.
Бұл модель негізгі көрсеткіштер бойынша XGBoost-пен салыстыруға болатын
нәтижелерді көрсетті, бірақ ол біршама жоғары есептеу шығындарымен
сипатталады, бұл оны шектеулі өндірістік желілер жағдайында қолдануды
шектейді.

Үшінші орында тұрған Random Forest моделі қалыпты есептеу талаптарында
жақсы тұрақтылық пен салыстырмалы түрде жоғары дәлдікті көрсетеді.
Дегенмен, оның күрделі сызықтық емес тәуелділіктерді модельдеу қабілеті
қатенің біршама жоғары мәндерінде көрінетін күшейту алгоритмдерінен
төмен.

Рейтингте төртінші орынды SVM моделі алады, ол ең төменгі Accuracy
мәндерін және нәтижелердің ең үлкен таралуын көрсетті. Бұл күрделі
техногендік қалдықтарды гранулометриялық талдаудың көп класты
мәселелерін шешуде осы әдістің шектеулі тиімділігін көрсетеді.

Алынған нәтижелер техногендік қалдықтардың гранулометриялық құрамын
автоматтандырылған анықтау үшін машиналық оқыту әдістерін қолданудың
тиімділігін растайды, алайда ұсынылған тәсілдің одан әрі дамуы
өнеркәсіптік өндіріс жағдайында оның практикалық қолданылуын кеңейтуге
бағытталған бірқатар перспективалық бағыттармен байланысты.

Әрі қарай зерттеудің негізгі бағыттарының бірі-офлайн талдаудан онлайн
түсірілімге көшу. Өнеркәсіптік камераларды пайдалана отырып, суреттерді
ағынмен алуды іске асыру қайта өңдеу процесінде материалдардың
түйіршіктелген құрамын үздіксіз мониторингтеуге мүмкіндік береді. Бұл
тұрғыда түсірілімді материалдың қозғалысымен синхрондау, діріл мен
біркелкі емес жарықтандырудың орнын толтыру, сондай-ақ сегменттеу және
анықтау алгоритмдерін динамикалық өзгеретін түсіру жағдайларына бейімдеу
өзекті міндеттер болып табылады.

Екінші маңызды бағыт әзірленген әдісті Құрылыс және тау-кен металлургия
салалары кәсіпорындарының өндірістік желілерімен интеграциялау болып
табылады. Технологиялық процестерді басқару контурына Машиналық оқыту
алгоритмдерін енгізу гранулометриялық құрамның ағымдағы сипаттамалары
негізінде Шикізатты ұсақтау, сұрыптау және мөлшерлеу режимдерін жедел
түзету мүмкіндіктерін ашады. Ол автоматтандыру жүйелерімен (SCADA, MES)
өзара әрекеттесу интерфейстерін әзірлеуді, сондай-ақ модельдердің нақты
уақыт режимінде тұрақты жұмыс істеуін қамтамасыз етуді талап етеді.

Үшінші перспективалық бағыт металлургиялық қождар, күл-қож материалдары,
байыту қалдықтары және құрылыс бөлшектеу қалдықтары сияқты техногендік
қалдықтардың әртүрлі түрлерін қосу арқылы датасеттің кеңеюі мен
әртараптандырылуына байланысты. Деректердің көлемі мен әртүрлілігін
арттыру модельдердің жалпылау қабілетін арттыруға, жеке материалдардың
ерекшеліктеріне сезімталдықты төмендетуге және әртүрлі кәсіпорындар мен
технологиялық схемалар арасындағы әдістің тасымалдануын қамтамасыз етуге
мүмкіндік береді.

Аталған бағыттарды іске асыру қайталама шикізат сапасын бақылаудың
зияткерлік жүйелерін қалыптастыру үшін алғышарттар жасайды және құрылыс
саласында цифрлық технологиялардың одан әрі дамуына ықпал етеді.

{\bfseries Қорытынды.} Зерттеу барысында кескінді алдын-ала өңдеудің, U-net
моделін қолдана отырып бөлшектерді сегментациялаудың, yolov5 негізіндегі
объектілерді анықтаудың, геометриялық белгілерді алудың және машиналық
оқыту алгоритмдерінің ансамблін (Random Forest, XGBoost және SVM)
қолдана отырып фракцияларды жіктеудің дәйекті кезеңдерін қамтитын қож
және күл түйіршіктерінің SEM кескіндерін талдауға негізделген кешенді
әдіс ұсынылды). Нәтижелердің дұрыстығы мен қайталануын қамтамасыз ету
үшін 5 есе кросс-валидацияны, сенімділік интервалдарын талдауды және
алынған үлестірімдерді МЕМСТ 5954.2-2020 және ASTM c136 стандарты
бойынша эталондық Елек талдауының нәтижелерімен салыстыруды қамтитын
верификация рәсімі іске асырылды.

Жұмыстың ғылыми жаңалығы - "компьютерлік көру-Машиналық оқыту" гибридті
архитектурасын техногендік қалдықтардың гранулометриялық құрамын анықтау
міндетіне қолдану, сондай-ақ тұрақтылықты, қайталануды және инженерлік
қолдануды ескере отырып, әртүрлі модельдерді кешенді статистикалық
салыстыру. Зертханалық үлгілер мен табиғи минералдарға бағытталған
көптеген зерттеулерден айырмашылығы, ұсынылған тәсіл жоғары
морфологиялық гетерогенділік пен бөлшектердің агломерациясы бар
материалдарға бейімделген.

Сандық нәтижелер дәлдіктің ең жақсы көрсеткіштерін XGBoost моделі (95\%
дейін Accuracy, MAE шамамен 6-7\%) көрсететінін көрсетті, u-net + yolov5
гибридті тәсілі бөлшектерді сегменттеу мен анықтаудың жоғары дұрыстығын
қамтамасыз етеді. Қате матрицалары мен кросс-валидация нәтижелерін
талдау ансамбльдік модельдердің тұрақтылығын және жіктеудің өрескел
қателіктерінің жоқтығын растады. Сапалық талдау кескін сапасына,
бөлшектердің жиналу тығыздығына және өте ұсақ фракцияларды талдауға
байланысты әдістің негізгі шектеулерін анықтады, бірақ бұл шектеулер
интегралды гранулометриялық үлестірудің айтарлықтай бұрмалануына
әкелмейтінін көрсетті.

Жұмыстың практикалық маңыздылығы құрылыс саласында пайдалану кезінде
техногендік қалдықтардың гранулометриялық құрамын жедел және
автоматтандырылған бақылау үшін ұсынылған әдісті қолдану мүмкіндігі
болып табылады. Бұл тәсілді іске асыру талдау уақытын бірнеше есе
қысқартуға, адам факторының әсерін барынша азайтуға және Шикізат сапасын
бақылаудың ауқымдылығын қамтамасыз етуге мүмкіндік береді.

Әрі қарайғы зерттеулердің болашағы онлайн Түсірілім мен ағынды кескінді
өңдеуге көшумен, дамыған алгоритмдерді өндірістік желілермен және
автоматтандыру жүйелерімен біріктірумен және техногендік қалдықтардың
әртүрлі түрлерімен датасетті кеңейтумен байланысты. Аталған бағыттарды
іске асыру қайталама материалдардың сапасын басқарудың Зияткерлік
жүйелерін қалыптастыру үшін негіз жасайды және құрылыс және қайта өңдеу
өнеркәсібінде цифрлық технологиялардың дамуына ықпал етеді.

Техногендік қалдықтардың гранулометриялық құрамын автоматтандырылған
түрде анықтаудың ұсынылған әдісін енгізу зертханалық бақылауға кететін
шығындарды қысқарту, технологиялық шешімдер қабылдаудың жеделдігін
арттыру және Шикізат шығынын азайту есебінен айқын экономикалық әсерді
қамтамасыз етеді. Айтарлықтай еңбек және уақыт ресурстарын қажет ететін
дәстүрлі елеуіш талдауларымен салыстырғанда, машиналық көру және
машиналық оқыту әдістерін қолдану талдау ұзақтығын 3-5 есе қысқартуға
мүмкіндік береді, бұл әсіресе материалдарды ағынмен өңдеу кезінде
маңызды.

Қолмен жұмыс істеу үлесінің төмендеуі және оператордың субъективті
факторын алып тастау зертханалардың пайдалану шығындарының төмендеуіне
және нәтижелердің қайталануының жоғарылауына әкеледі. Гранулометриялық
құрамды автоматтандырылған талдау шикізаттың фракциялық құрамын дәлірек
бақылауға мүмкіндік береді, бұл өз кезегінде құрылыс қоспаларын өндіруде
тұтқыр компоненттер мен түзету қоспаларының артық шығынын азайтады. Дәл
PSD деректеріне негізделген формулаларды оңтайландыру материалдың түріне
және өндіріс технологиясына байланысты шикізат ресурстарын 5-10\%
үнемдеуге мүмкіндік береді деп есептеледі.

\emph{{\bfseries Қаржыландыру.} Мақалада сипатталған зерттеу "AG TECH" ЖШС
қаржыландыратын "Тау-кен металлургия өндірісінің барлық кезеңдеріндегі
өнім сапасын басқару жүйесінің цифрлық қосарлануы" бастамалық жобасы
аясында жүргізілді.}
\end{multicols}

\begin{center}
{\bfseries Әдебиеттер}
\end{center}

\begin{refs}
1. Baraian A., et al. Computing Particle Size Distribution of Mineral
Rocks using Deep Learning-based Instance Segmentation// 2022 10th
European Workshop on Visual Information Processing (EUVIP). -2022. -P.
1-6. \href{https://doi.org/10.1109/EUVIP53989.2022.9922748}{DOI
10.1109/EUVIP53989.2022.9922748}.

2. Nie H., et al. Image Segmentation Method on Quartz Particle‑Size
Detection by Deep Learning Network// Minerals. -2023. -Vol.12(12):1479.
\href{https://doi.org/10.3390/min12121479}{DOI 10.3390/min12121479}.

3. Latif A., et al. Deep‑Learning-Based Automatic Mineral Grain
Segmentation and Recognition// Minerals. -2022. -Vol.12(4):455.
\href{https://doi.org/10.3390/min12040455}{DOI 10.3390/min12040455}.

4. Neubauer Т., et al. Multiscale attention-based instance segmentation
for measuring crystals with large size variation // Computer Vision and
Pattern Recognition (cs.CV). -2024. -arXiv:2401.0339. DOI
\href{https://doi.org/10.1109/TIM.2023.3345916}{10.1109/TIM.2023.3345916}.

5. Jiang J., et al. Res-UNet Ensemble Learning for Semantic Segmentation
of Mineral Optical Microscopy Images //Minerals. -2024.
-Vol.14(12):1281. \href{https://doi.org/10.3390/min14121281}{DOI
10.3390/min14121281}.

6. Zhou X., He Y. Deep Ensemble Learning-Based Sensor for Flotation
Froth Image Recognition// Sensors. -2024. -Vol.~24(15):5048.
\href{https://doi.org/10.3390/s24155048}{DOI 10.3390/s24155048}.

7. Liu Y., et al. Deep learning in image segmentation for mineral
production: A review// Computers \& Geosciences. -2023.
-\href{https://www.sciencedirect.com/journal/computers-and-geosciences/vol/180/suppl/C}{Vol.
180}:105455. \href{https://doi.org/10.1016/j.cageo.2023.105455}{DOI
10.1016/j.cageo.2023.105455}.

8. Luo G. Online detection system for ore PSD based on deep
learning//~Journal of Computer, Signal, and System Research. - 2025.
-Vol.2(3). - P.1- 16. \href{https://doi.org/10.71222/r6qws842}{DOI
10.71222/r6qws842}.

9. Tae Soo E., et al. Development and application of a user-friendly
general-purpose predictive simulation tool for two-dimensional flow
analysis//Environmental Modelling \& Software. -2023. -Vol.163:105665.
\href{https://doi.org/10.1016/j.envsoft.2023.105665}{DOI
10.1016/j.envsoft.2023.105665}.

10. Gotkowski K., et al. ParticleSeg3D: A Scalable Out-of-the-Box Deep
Learning Segmentation Solution for Individual Particle Characterization
from Micro CT Images in Mineral Processing and Recycling//
Computer Vision and Pattern Recognition (cs.CV). -2023.DOI
10.48550/arXiv.2301.13319.

11. Martins L., et al. Deep Learning Domain Adaptation for Minerals
Semantic Segmentation in Reflected Light microscopy images// Preprints.
-2024.- P.2-28. DOI 10.20944/preprints202412.0572.v1.

12. Qiu Y., et al. Minima-YOLO: A Lightweight Identification Method for
Lithium Mineral Components Under a Microscope Based on YOLOv8// Sensors.
-2025. --Vol.25(7):2048. \href{https://doi.org/10.3390/s25072048}{DOI
10.3390/s25072048}.

13. Yimin L., et al. Flow characterization of compressible biomass
particles using multiscale experiments and a hypoplastic model// Powder
Technology. -2021.-Vol.383. -P.396-409.
\href{https://doi.org/10.1016/j.powtec.2021.01.027}{DOI
10.1016/j.powtec.2021.01.027}.

14. Kim B., et al. Multi-scale flume investigation of the influence of
cylindrical baffles on the mobility of landslide debris//Engineering
Geology. - 2023.-Vol.314:107012. DOI
\href{https://doi.org/10.1016/j.enggeo.2023.107012}{10.1016/j.enggeo.2023.107012}.

15. Акишев К., Арынгазин К.Ш., Карпов В.И. Применение методов
кластерного анализа для статистической оценки качества металлургического
шлака Павлодарского филиала ТОО «Кастинг»//Проблемы автоматики и
управления. -2019. -№ 2. - С.79-87.

16. Matthew J., et al. An industrial 3D vision system for size
measurement of iron ore green pellets using morphological image
segmentation//
\href{https://www.sciencedirect.com/journal/minerals-engineering}{Minerals
Engineering}. -2008. -\href{Vol.\%2021(5)}{Vol.21(5)}. -P.405-415.
\href{https://doi.org/10.1016/j.mineng.2007.10.020}{DOI
10.1016/j.mineng.2007.10.020}.

17. Andersson T., et al. A machine vision system for estimation of size
distributions by weight of limestone
particles//\href{https://www.sciencedirect.com/journal/minerals-engineering}{Minerals
Engineering}.-2012.-\href{Vol.\%2025(1}{Vol.25(1}).-P.38-46.
\href{https://doi.org/10.1016/j.mineng.2011.10.001}{DOI
10.1016/j.mineng.2011.10.001}.

18. Kemeny J. M., Devgan A., Hagaman R. M., Wu X. Analysis of Rock
Fragmentation Using Digital Image Processing //Journal of Geotechnical
Engineering (ASCE). -1993. - Vol.119(7). -P.1144-1160. DOI
\href{https://doi.org/10.1061/(ASCE)0733-9410(1993)119:7(1144)}{10.1061/(ASCE)0733-9410(1993)119:7(1144}).

19. Andersson T., Thurley M. J. Visibility Classification of Rocks in
Piles//IEEE Computer Societ\emph{. -}2008. -P.207--213.
\href{https://doi.org/10.1109/DICTA.2008.85}{DOI 10.1109/DICTA.2008.85}.

20. ГОСТ 5954.2-2020. Ситовый анализ класса крупности не менее 20 мм.
//Стандартинформ. -2020. - 11с.

21. ASTM C136 // C136M -19. Standard Test Method for Sieve Analysis of
Fine and Coarse Aggregates//Book of standarts. -2025\emph{.}
-Vol.04.02\emph{.} DOI~10.1520/C0136\_C0136M-19.
\end{refs}

\begin{center}
{\bfseries References}
\end{center}

\begin{refs}
1. Baraian A., et al. Computing Particle Size Distribution of Mineral
Rocks using Deep Learning-based Instance Segmentation// 2022 10th
European Workshop on Visual Information Processing (EUVIP). -2022. -P.
1-6. \href{https://doi.org/10.1109/EUVIP53989.2022.9922748}{DOI
10.1109/EUVIP53989.2022.9922748}.

2. Nie H., et al. Image Segmentation Method on Quartz Particle‑Size
Detection by Deep Learning Network// Minerals. -2023. -Vol.12(12):1479.
\href{https://doi.org/10.3390/min12121479}{DOI 10.3390/min12121479}.

3. Latif A., et al. Deep‑Learning-Based Automatic Mineral Grain
Segmentation and Recognition// Minerals. -2022. -Vol.12(4):455.
\href{https://doi.org/10.3390/min12040455}{DOI 10.3390/min12040455}.

4. Neubauer Т., et al. Multiscale attention-based instance segmentation
for measuring crystals with large size variation // Computer Vision and
Pattern Recognition (cs.CV). -2024. -arXiv:2401.0339. DOI
\href{https://doi.org/10.1109/TIM.2023.3345916}{10.1109/TIM.2023.3345916}.

5. Jiang J., et al. Res-UNet Ensemble Learning for Semantic Segmentation
of Mineral Optical Microscopy Images //Minerals. -2024.
-Vol.14(12):1281. \href{https://doi.org/10.3390/min14121281}{DOI
10.3390/min14121281}.

6. Zhou X., He Y. Deep Ensemble Learning-Based Sensor for Flotation
Froth Image Recognition// Sensors. -2024. -Vol.~24(15):5048.
\href{https://doi.org/10.3390/s24155048}{DOI 10.3390/s24155048}.

7. Liu Y., et al. Deep learning in image segmentation for mineral
production: A review// Computers \& Geosciences. -2023.
-\href{https://www.sciencedirect.com/journal/computers-and-geosciences/vol/180/suppl/C}{Vol.
180}:105455. \href{https://doi.org/10.1016/j.cageo.2023.105455}{DOI
10.1016/j.cageo.2023.105455}.

8. Luo G. Online detection system for ore PSD based on deep
learning//~Journal of Computer, Signal, and System Research. - 2025.
-Vol.2(3). - P.1- 16. \href{https://doi.org/10.71222/r6qws842}{DOI
10.71222/r6qws842}.

9. Tae Soo E., et al. Development and application of a user-friendly
general-purpose predictive simulation tool for two-dimensional flow
analysis//Environmental Modelling \& Software. -2023. -Vol.163:105665.
\href{https://doi.org/10.1016/j.envsoft.2023.105665}{DOI
10.1016/j.envsoft.2023.105665}.

10. Gotkowski K., et al. ParticleSeg3D: A Scalable Out-of-the-Box Deep
Learning Segmentation Solution for Individual Particle Characterization
from Micro CT Images in Mineral Processing and Recycling//
Computer Vision and Pattern Recognition (cs.CV). -2023.DOI
10.48550/arXiv.2301.13319.

11. Martins L., et al. Deep Learning Domain Adaptation for Minerals
Semantic Segmentation in Reflected Light microscopy images// Preprints.
-2024.- P.2-28. DOI 10.20944/preprints202412.0572.v1.

12. Qiu Y., et al. Minima-YOLO: A Lightweight Identification Method for
Lithium Mineral Components Under a Microscope Based on YOLOv8// Sensors.
-2025. --Vol.25(7):2048. \href{https://doi.org/10.3390/s25072048}{DOI
10.3390/s25072048}.

13. Yimin L., et al. Flow characterization of compressible biomass
particles using multiscale experiments and a hypoplastic model// Powder
Technology. -2021.-Vol.383. -P.396-409.
\href{https://doi.org/10.1016/j.powtec.2021.01.027}{DOI
10.1016/j.powtec.2021.01.027}.

14. Kim B., et al. Multi-scale flume investigation of the influence of
cylindrical baffles on the mobility of landslide debris//Engineering
Geology. - 2023.-Vol.314:107012. DOI
\href{https://doi.org/10.1016/j.enggeo.2023.107012}{10.1016/j.enggeo.2023.107012}.

15. Akishev K., Aryngazin K.Sh., Karpov V.I. Primenenie metodov
klasternogo analiza dlja statisticheskoj ocenki kachestva
metallurgicheskogo shlaka Pavlodarskogo filiala TOO «Kasting»//Problemy
avtomatiki i upravlenija. -2019.-№ 2. - S.79-87 {[}in Russian{]}.

16. Matthew J., et al. An industrial 3D vision system for size
measurement of iron ore green pellets using morphological image
segmentation//
\href{https://www.sciencedirect.com/journal/minerals-engineering}{Minerals
Engineering}. -2008. -\href{Vol.\%2021(5)}{Vol.21(5)}. -P.405-415.
\href{https://doi.org/10.1016/j.mineng.2007.10.020}{DOI
10.1016/j.mineng.2007.10.020}.

17. Andersson T., et al. A machine vision system for estimation of size
distributions by weight of limestone
particles//\href{https://www.sciencedirect.com/journal/minerals-engineering}{Minerals
Engineering}.-2012.-\href{Vol.\%2025(1}{Vol.25(1}).-P.38-46.
\href{https://doi.org/10.1016/j.mineng.2011.10.001}{DOI
10.1016/j.mineng.2011.10.001}.

18. Kemeny J. M., Devgan A., Hagaman R. M., Wu X. Analysis of Rock
Fragmentation Using Digital Image Processing //Journal of Geotechnical
Engineering (ASCE). -1993. - Vol.119(7). -P.1144-1160. DOI
\href{https://doi.org/10.1061/(ASCE)0733-9410(1993)119:7(1144)}{10.1061/(ASCE)0733-9410(1993)119:7(1144}).

19. Andersson T., Thurley M. J. Visibility Classification of Rocks in
Piles//IEEE Computer Societ\emph{. -}2008. -P.207--213.
\href{https://doi.org/10.1109/DICTA.2008.85}{DOI 10.1109/DICTA.2008.85}.

20. GOST 5954.2-2020. Sitovyj analiz klassa krupnosti ne menee 20 mm.
//Standartinform. -2020. -11s. {[}in Russian{]}

21. ASTM C136 // C136M -19. Standard Test Method for Sieve Analysis of
Fine and Coarse Aggregates//Book of standarts. -2025\emph{.}
-Vol.04.02\emph{.} DOI~10.1520/C0136\_C0136M-19.
\end{refs}

\begin{info}
\hspace{1em}\emph{{\bfseries Авторлар туралы мәлімет}}

Акишев К.М. - т. ғ. к., қауымдастырылған профессор (доцент), Қ.Құлажанов
атындағы Қазақ технология және бизнес университеті, Астана, Қазақстан,
e-mail: akmail04cx@mail.ru;

Акишева Л. - Назарбаев зияткерлік мектебі, Астана, Қазақстан, e-mail:
L.Ak@mail.ru;

Нұртай Ж.Т. - PhD докторы, асс. профессор, Қ.Құлажанов атындағы Қазақ
технология және бизнес университеті, Астана, Қазақстан, e-mail:
zhadira\_nurtai@mail.ru.

\hspace{1em}\emph{{\bfseries Information about the authors}}

Akishev K. M. - Candidate of Technical Sciences, Associate Professor
(Associate Professor), K. Kulazhanov Kazakh University of Technology and
Business, Astana, Kazakhstan,e-mail:akmail04cx@mail.ru;

Akisheva L.- Nazarbayev Intellectual School, Astana, Kazakhstan, e-mail:
L. Ak@mail.ru;

Nurtai Zh.T. - PhD, Associate Professor, K. Kulazhanov Kazakh University
of Technology and Business, Astana, Kazakhstan, e-mail:
zhadira\_nurtai@mail.ru.
\end{info}
